\section{The ordinary-exit criterion and the two reconstructions}
\label{sec:criterion}
Let $D\subset X$ be a reduced normal-crossings divisor and $U=X\setminus D$.
Write $\mathcal E=\Exit(X,D)$, with arrows pointing from a deeper stratum
to a more open one. At a point $x$ having local branch set $I_x$, put
$\Lambda_x=\Z^{I_x}$ and let $\mu_x:\Lambda_x\to\pi_1(U)$ be the
homomorphism of positive normal meridians, with a chosen connecting path.
Every assertion is componentwise. Neither $\mathcal E$ nor its classifying
space is assumed to be a one-type.

\subsection{Normal links, flags, and specialization}

\begin{theorem}[The normal-link mapping-space formula]\label{P1a1}
\status{Proved}.
Let $S,T$ be distinct incident strata, of depths $d>e$.
There is a finite incidence cover $F_{ST}\to S$ selecting the local
branches that remain zero in $T$, and a torus bundle
\[
 (S^1)^{d-e}\longrightarrow L_{ST}\longrightarrow F_{ST}.
\]
Its endpoint maps give a coherently defined map $L_{ST}\to S\times T$.
For $s\in S,t\in T$,
\[
 \Map_{\mathcal E}(s,t)\simeq
 \hofib_{(s,t)}(L_{ST}\longrightarrow S\times T). \tag{1}
\]
Within one stratum the mapping space is its path space; in particular
$\Map_{\mathcal E}(s,s)\simeq\Omega_sS$.
Nonincident strata have empty crossing mapping spaces.
These assertions hold on stack homotopy types, retaining all inertia.
\end{theorem}
\begin{proof}
In a coordinate chart write $D$ as $z_1\cdots z_d=0$ at $S$.
A local generization to $T$ selects the $e$ coordinates that stay zero.
In the other $d-e$ coordinates a short exit has positive radius and a
nonzero angular coordinate. The radial choices are contractible and
the angular choices form $(S^1)^{d-e}$. Continuation along $S$
permutes branch selections and glues angles by the normal line-bundle
transitions. This produces the incidence cover and torus bundle,
rather than an unjustified globally labelled product.

The holink consists of paths starting in $S$ and immediately entering
$T$. Normal-coordinate contraction identifies its homotopy type with
the angular model, over the starting point and up to homotopy over
the endpoint. Allow arbitrary paths in $T$ to the prescribed endpoint
and paths in $S$ to the prescribed starting point. This is exactly
the homotopy fibre in (1). The holink-to-exit comparison is
\citet[Theorem 4.9]{Miller2009}; normal-crossing charts satisfy its
local tameness hypotheses. For the local-to-global use of its finite
stratification formulation, a compact parameter family of paths is
covered by finitely many such charts and their finitely many local
incidence types. The conical exit construction and
exodromy are those of \citet[Appendix A.9]{LurieHA2017}.

Within a single stratum no direction is exited, so one obtains its
fundamental infinity-groupoid. An exit cannot meet a stratum outside
the incidence closure order. Finally apply these constructions on
every level of an étale atlas and its entire Čech nerve. Branch
covers, links, and flag comparisons commute with étale pullback.
Descent of that coherent diagram proves the stack assertion, as in
\citet[Proposition 3.16 and Corollary 3.17]{JansenToolbox2024}.
No coarse quotient is substituted.
\end{proof}

\begin{theorem}[The normal-bundle and coherent-flag model]\label{P1a2}
\status{Proved}.
On a branch-labelled open incidence stratum $D_J^\circ$, let $N_j$
be its normal lines and put
\[
 P_J^\times=\prod_{j\in J,D_J^\circ}N_j^\times,\qquad
 P_J=P_J^\times/(\R_{>0})^J .
\]
The punctured local tube has homotopy type $P_J^\times\simeq P_J$.
These are the angular interiors of the corresponding closed faces
of $X^{\log}$. Specialization and its coherences are supplied by
closed-face inclusions and all incidence flags.
Let $\mathcal E_D=\Exit(D)$, $\mathcal E_U=\Pi_\infty U$, and
\[
 \mathcal H_X=\Fun(\Delta^1,\mathcal E)
 \times_{\mathcal E\times\mathcal E}
             (\mathcal E_D\times\mathcal E_U). \tag{2}
\]
Its objects are crossing exits and its morphisms are coherent squares.
Then
\[
                  B\mathcal H_X\simeq\partial X^{\log}. \tag{3}
\]
Strict positive-rank flags give a semisimplicial resolution. A full
simplicial resolution also includes repeated ranks and identity
degeneracies. The manuscript's fundamental-groupoid normal diagram
means applying $\Pi_1$ to every full flag space with its coherent
maps; it does not mean forming the crossing category of $h\mathcal E$.
\end{theorem}
\begin{proof}
The quotient $P_J^\times\to P_J$ is a principal bundle with
contractible fibre $(\R_{>0})^J$, and hence is a weak equivalence
on classifying homotopy types. On each manifold chart one may
choose Hermitian metrics and a differentiable tubular neighborhood:
a nonzero vector has a positive length and a unit direction, and
the radial factor contracts. This gives the polar-coordinate
description of punctured coordinate discs. For a general étale
stack these are atlas-level choices, not an invariant metric or
an equivariant tubular embedding.

The oriented blow-up replaces a vanishing coordinate by
$(r_j,\theta_j)\in\R_{\ge0}\times S^1$. A chosen branch
intersection defines a closed face. Its angular interior is the
torus bundle above. The interior inclusion is a homotopy
equivalence on each manifold level of the étale Čech nerve:
compatible collars push the remaining boundary coordinates
slightly positive. Over noncompact strata choose positive collar
functions locally finitely. The inclusion maps themselves are
natural, so these levelwise equivalences descend without choosing
equivariant collars. Forgetting branch
conditions is an actual face inclusion. Transporting it across
these equivalences gives the specialization maps; chains of
inclusions give their homotopies and all higher compatibilities.
Thus inverse equivalences are not chosen as unrelated,
incoherent tubular embeddings.

Alternatively use deformation to the normal cone. The nonzero
deformation fibres are the original tubular pair and the zero
fibre is the normal pair. On punctured normal coordinates the
intervening positive radial parameters are contractible.
Apply this to the entire flag system, then descend through
branch covers and the étale Čech nerve. This gives the
homotopy-level specialization despite the analytic transition
or transvection at scan 146. No global equivariant tubular
embedding, invariant metric, or holomorphic normal-bundle
linearization is asserted. The logarithmic comparison
is compatible with \citet{KatoNakayama1999}.

In a conical chart a crossing path is determined, up to
contractible radial choices, by its angular direction and
its endpoint paths. Groupoid completion of the endpoint
actions leaves the angular boundary. A chain of squares in
(2) records a chain of incident faces with angular data.
Closed-face inclusions have collars and are cofibrant, so
their homotopy-colimit realization is their union.
The usual Čech nerve is indexed by subsets of faces;
its barycentric flag refinement is the strict-flag
semisimplicial model. Repetitions provide degeneracies.
The chartwise equivalences glue and descend, proving (3).

A loop in a crossing mapping space can be lost in
$h\mathcal E$ although it survives in $\Pi_1$ of the full
link. The examples below exhibit this distinction.
\end{proof}

\subsection{The exact criterion and the failure of adjacent tests}

\begin{theorem}[The ordinary-exit criterion]\label{P1b1}
\status{Proved}. The following are equivalent:
\begin{enumerate}
\item $\mathcal E$ is equivalent to an ordinary category.
\item Every open incidence stratum is aspherical and every full
      $\mu_x:\Lambda_x\to\pi_1(U)$ is injective.
\item Every open incidence stratum is aspherical and every relative
      normal-meridian lattice for every incident pair injects into
      the group of the more open stratum.
\end{enumerate}
Full injection implies all relative injections even without
asphericity. Testing only immediately adjacent strata is insufficient.
\end{theorem}
\begin{proof}
An infinity-category is equivalent to an ordinary category
exactly when its mapping spaces are homotopy-discrete.
Here the endomorphism spaces are $\Omega S$. Their
homotopy-discreteness is equivalent to $\pi_i(S)=0$ for $i\ge2$.

Assume the strata are aspherical. Each finite incidence cover
is aspherical. The torus-bundle homotopy sequence shows that
$L_{ST}$ is aspherical and that its fibre lattice injects.
The kernel of
\[
 \pi_1(L_{ST})\longrightarrow\pi_1(S)\times\pi_1(T)
\]
is exactly the kernel of that fibre lattice in $\pi_1(T)$:
the incidence cover injects on $\pi_1$ into $S$, so an
element in the displayed kernel lies in the fibre lattice.
For a map between aspherical spaces, every component of
its homotopy fibre is a $K(\ker,1)$. Formula (1) is
therefore homotopy-discrete exactly when the relative
kernel vanishes. This proves the third equivalence and
the necessity of the full maps by taking $T=U$.

For the reduction from full to relative maps, write the
source branch set as $I=J\sqcup K$, with $J$ still
vanishing in $T$. Push the full local torus into the
normal torus bundle $P_T\to T$. If $v\in\Z^K$ dies
in $\pi_1(T)$, exactness says its lift in $\pi_1(P_T)$
equals the image of some $w\in\Z^J$. This does not
require $\Z^J\to\pi_1(P_T)$ to inject.
The full vector $v-w$ consequently dies in $\pi_1(U)$.
Full injection implies $v-w=0$, and disjoint supports
give $v=w=0$. No asphericity was used in this reduction.
It is compatible with incidence covers and stack basepoints.
The next example refutes the adjacent shortcut.
\end{proof}

\begin{counterexample}[Adjacent injections with a full kernel]\label{P1b2}
\status{Refuted}.
Even when every stratum is aspherical, injectivity for all adjacent
pairs does not imply goodness.
\end{counterexample}
\begin{proof}
Start with the axes $A=\{x=0\},B=\{y=0\}$ in $\C^2$.
Blow up $p=(0,1)\in A$, write $E$ for the exceptional
curve, and remove a point $r\in E$ different from
$q=E\cap A'$. Use the reduced total boundary.
The open complement remains $U=(\C^*)^2$, and the
boundary curve strata are
\[
 A^\circ=\C\setminus\{0,1\},\quad B^\circ=\C^*,\quad
 E^\circ=\PP^1\setminus\{q,r\}=\C^* .
\]
There are two point strata. All strata are aspherical.
The curve-to-open meridians are $e_x,e_y,e_x$ in $\Z^2$,
so each is injective. At either node, each relative
node-to-curve meridian is a nontrivial peripheral loop
in one of these punctured curves. Its free group is
torsion-free, so the adjacent map injects.
At $q$, however, the full map is
\[
 \Z^2\longrightarrow\Z^2,\qquad (a,b)\longmapsto(a+b)e_x,
\]
with kernel $\Z(1,-1)$. Individual cyclic injection
cannot replace independence of the full lattice.
The integer check is in \texttt{criterion\_checks.py}.
\end{proof}

\begin{proposition}[The missing normal transgression]\label{P1b3}
\status{Proved}.
For a rank-$r$ normal torus bundle over a branch-labelled stratum,
\[
 \pi_2(S)\xrightarrow{\partial}\Z^r
 \longrightarrow\pi_1(P_S)\longrightarrow\pi_1(S)\longrightarrow1.
 \tag{4}
\]
The lattice injects exactly when $\partial=0$.
Its coordinates are the first Chern numbers of the normal
lines on spherical classes. Thus $\pi_2(S)=0$ is sufficient
but not necessary. The lattice quotient may have torsion.
In (37), the base must be the appropriate open incidence
stratum; substituting the closed intersection changes the claim.
\end{proposition}
\begin{proof}
This is the homotopy exact sequence of the torus fibration.
Pull a normal line back to a representative sphere and
trivialize over its hemispheres. Its equatorial clutching
function has degree its first Chern number. Lifting that
equator gives the same integer as the connecting map.
Taking the product gives all coordinates of $\partial$.

For $\mathcal O_{\PP^1}(1)$ the connecting map is
multiplication by one. For a trivial line over $\PP^1$
it is zero although $\pi_2(\PP^1)=\Z$.
A conic in $\PP^2$ has normal degree four, so its
circle-bundle group is $\Z/4$, exhibiting torsion.

For the coordinate axes, the first component's open
stratum is $\C^*$. Its punctured normal tube inside
the open complement has group $\Z^2$, mapping to
$\pi_1(\C^*)=\Z$ with rank-one kernel.
Replacing the base by the closed axis $\C$ gives
zero target and rank-two kernel. The proposed
sequence cannot describe both bases. This supplies
the correction to (37) and its later reminder in
MPL 149, scan 134.
\end{proof}

\subsection{Three operations which must be distinguished}

\begin{theorem}[Derived reconstruction and ordinary truncation]\label{P1c1}
\status{Proved}.
Let $\mathcal T$ be the complete normal-face/link diagram,
refined by open incidence strata and all flags.
\begin{enumerate}
\item Without goodness, its full homotopy types reconstruct $X$
      and its constructible infinity-topos. Its fundamental-groupoid
      diagram always reconstructs $\Pi_1X$, and $h\mathcal E$
      classifies constructible sets.
\item If all open strata are aspherical, the refined values can
      be replaced objectwise by $B\Pi_1$ without changing their
      homotopy colimit. The full ambient homotopy type is recovered
      from this derived groupoid diagram under asphericity alone.
\item With full meridian injection as well,
      $\mathcal E\simeq h\mathcal E$. Both the genuine constructible
      infinity-topos and the infinity-topos associated with
      ordinary constructible sets then have shape $X$.
\end{enumerate}
Item two concerns the refined open-stratum/link diagram;
it does not assert asphericity of closed branch intersections
or identify its homotopy colimit with $B(h\mathcal E)$.
\end{theorem}
\begin{proof}
A normal-crossing collar glues open strata to their normal
directions. Chartwise this is gluing coordinate cones
and faces, which becomes a homotopy pushout using collars.
Iterate through the finite depth and retain all flags.
An open Čech cover and an étale atlas give the global
and stack homotopy-colimit reconstruction.
Equivalently exodromy identifies the constructible
space-valued sheaves with $\Fun(\mathcal E,\mathrm{Spaces})$,
and groupoid completion gives the homotopy type of $X$;
see \citet[Appendix A.9]{LurieHA2017} and
\citet[Proposition 3.24 and Corollary 3.23]{JansenToolbox2024}.
The collar argument supplies the diagram used here.

The fundamental-groupoid functor is left adjoint to
the inclusion of one-types and preserves colimits
formed in groupoids. No asphericity is needed for
the $\Pi_1$ conclusion. Every set-valued functor
on an infinity-category factors uniquely through
its homotopy category. Exodromy therefore also gives
$\Fun(h\mathcal E,\mathrm{Set})$ without goodness.

Under asphericity, refined diagram values are
finite covers of open strata, iterated normal-torus
bundles, and contractible radial parameters.
Their homotopy exact sequences prove asphericity.
The natural maps to $B\Pi_1$ are objectwise
equivalences of the full coherent diagram.
They remain equivalences after homotopy colimits.
Deleting triangular homotopies is a different
operation and can change even $\Pi_1$, as Part 5 shows.

Finally apply Theorem~\ref{P1b1}. For any small category
$\mathcal C$, global sections of a constant space
$K$ in its presheaf infinity-topos are
$\lim_{\mathcal C}K=\Map(B\mathcal C,K)$.
Thus its shape is $B\mathcal C$.
Apply this to $\mathcal E$ and to $h\mathcal E$.
Their shapes agree under goodness. The stronger
ordinary-topos consequence requires this extra
hypothesis and does not follow from item two alone.
\end{proof}

\begin{counterexample}[The separate roles of the two hypotheses]\label{P1c2}
\status{Refuted}.
Asphericity alone does not identify the ordinary constructible
shape with $X$. Full meridian injection alone does not recover
$X$ from its fundamental-groupoid diagram.
\end{counterexample}
\begin{proof}
For $(\PP^1,\{\infty\})$ both strata are contractible
and the link is $S^1$. The derived groupoid diagram
\[
                    *\longleftarrow B\Z\longrightarrow *
\]
has homotopy pushout $S^2$. But the meridian dies
in the open $\C$. The crossing exit space is $S^1$,
so $h\mathcal E=[1]$ and its classifying space is
contractible. The ordinary consequence fails,
although derived reconstruction succeeds.

For the other direction take
\[
             (X,D)=(\PP^1\times\C,\PP^1\times\{0\}).
\]
Its normal meridian maps isomorphically onto
$\pi_1(U)=\Z$, while the closed stratum is
not aspherical. Its closed-link-open groupoid
diagram is
\[
                 *\longleftarrow B\Z\xrightarrow{\id}B\Z ,
\]
with contractible homotopy pushout.
This is exactly the same diagram as for
$(\C^2,\C\times\{0\})$, although the two
ambient homotopy types are $S^2$ and a point.
Thus those groupoid data cannot distinguish them.
\end{proof}

\subsection{Shapes, fibres, and the first possible obstruction}

\begin{theorem}[The two shapes and their gluing data]\label{P1d1}
\status{Proved}.
The genuine constructible infinity-topos has shape
$B\mathcal E\simeq X$. The infinity-topos associated
with the ordinary constructible-set topos has shape
$B(h\mathcal E)$. The canonical comparison is
\[
            q:X\simeq B\mathcal E\longrightarrow B(h\mathcal E). \tag{5}
\]
With two strata, closed $S$ and open $T$, and link $L$,
the genuine shape is
\[
                              S\cup_L^hT. \tag{6}
\]
More strata require every incidence flag.
If $S,T,L$ are aspherical, with groups $A,B,G$
and endpoint homomorphism $G\to A\times B$,
the ordinary crossing homotopy type uses
$\im(G\to A\times B)$ in place of $G$.
\end{theorem}
\begin{proof}
The shape assertions follow from the constant
global-sections calculation above. The closed
stratum has a mapping-cylinder neighborhood with
boundary $L$; gluing to the open complement
gives (6). Shrinking collars does not change
the diagram over its endpoints.

The crossing mapping space is the homotopy
fibre of $BG\to B(A\times B)$. Its components
are cosets of the image, each with group
the kernel. For the ordinary exit category
only the component set remains, acted on
by $A\times B$. The homotopy quotient of
this transitive action is the classifying
space of its stabilizer, the image subgroup.
Before truncation that quotient is $BG$.
This proves the precise change in the
crossing data. More strata require the
same operation coherently over full flags.
\end{proof}

\begin{proposition}[Explicit comparison fibres]\label{P1d2}
\status{Proved}.
The comparison fibres in (5) for
$(\PP^1,\{\infty\})$, $(\PP^n,\PP^{n-1})$,
with $n\ge1$, and the product of two copies of
$(\PP^1,\{\infty\})$ are respectively
\[
                    S^2,\qquad\PP^n,\qquad S^2\times S^2 .
\]
Their first nonzero fibre groups are $\Z,\Z,\Z^2$
in degree two.
\end{proposition}
\begin{proof}
In the first case the ordinary category is $[1]$.
For the hyperplane pair the open is $\C^n$,
the closed stratum has trivial fundamental group,
and the link is the Hopf sphere $S^{2n-1}$.
Its map to $\PP^{n-1}$ is the circle bundle,
so (1) gives crossing mapping space $S^1$;
the closed endomorphism space is
$\Omega\PP^{n-1}$. Passing to components
again gives $[1]$, of contractible shape.
The comparison fibre is therefore $\PP^n$.

Exit categories of product stratifications
are products: paths, exit conditions, and
homotopies are coordinatewise. The third
ordinary category is $[1]\times[1]$,
again contractible, and its genuine shape
is $\PP^1\times\PP^1$.
The cell structures and product formula
give the stated second homotopy groups.
The cellular Betti-number checks are in
\texttt{criterion\_checks.py}.
\end{proof}

\begin{theorem}[Connectivity and the first comparison obstruction]\label{P1d3}
\status{Proved}.
For every small infinity-category $\mathcal E$,
$B\mathcal E\to B(h\mathcal E)$ is two-connected.
If each component of every mapping space is
$(r-1)$-connected, $r\ge1$, it is
$(r+1)$-connected.

Choose a well-based enriched-category model
with discrete objects, and put
\[
 V_p=\coprod_{x_0,\ldots,x_p}
       \Map(x_0,x_1)\times\cdots\times\Map(x_{p-1},x_p),
 \qquad W_p=\pi_0V_p .
\]
Their realizations are $B\mathcal E$ and $B(h\mathcal E)$.
Fix a component $Y$ of the latter. Use its fundamental-groupoid
local system of basepoint paths to form the covering bar
$\widetilde W_\bullet$, and pull this cover back to
$\widetilde V_\bullet$. Their realizations are the universal
cover of $Y$ and its pullback to $B\mathcal E$.
In particular the faces transport lifts along the corresponding
edges; this is not an unrelated levelwise choice of lifts.
For its homotopy fibre $F$,
\[
 \pi_{r+1}F\cong
 \coker\left(
 H_r(\widetilde V_2;\Z)
 \xrightarrow{d_0-d_1+d_2}
 H_r(\widetilde V_1;\Z)\right) . \tag{7}
\]
This is an isomorphism of $\Z[\pi_1Y]$-modules.
It is the first possible obstruction, not necessarily
a nonzero one. The first nonzero disagreement need not
occur in degree two, even for smooth algebraic NC pairs.
\end{theorem}
\begin{proof}
Use a cofibrant simplicial-category model,
so the displayed composition is strict and
the nerve degeneracies are cofibrations.
Each component of $V_p$ is a product of
$(r-1)$-connected spaces, and $V_p\to W_p$
is $r$-connected; at degree zero it is an
isomorphism. In relative realization,
a cell in simplicial degree $p$ acquires
$p$ extra dimensions. The first possible
relative cells in positive simplicial
degree occur in dimension $r+2$, at $p=1$.
The skeletal filtration of the simplicial
mapping cylinder therefore proves
$(r+1)$-connectivity. For $r=1$ this
applies to every mapping-space component.

Over the universal cover of $Y$ the relative
realization spectral sequence is
\[
 E^1_{p,j}=H_j(\widetilde W_p,\widetilde V_p;\Z)
 \Longrightarrow H_{p+j}(\widetilde Y,\widetilde X;\Z).
\]
The degree-zero column vanishes.
For $p>0$ the rows $j\le r$ vanish and
\[
 H_{r+1}(\widetilde W_p,\widetilde V_p)
       \cong H_r(\widetilde V_p).
\]
In total degree $r+2$ the only position
is $(1,r+1)$. Its incoming first
differential is $d_0-d_1+d_2$, and the
outgoing one lands in the zero column.
Higher incoming differentials start
in vanishing lower rows; higher
outgoing ones leave the first quadrant.
The relative homology group is thus
the cokernel in (7).
The already proved connectivity and
relative Hurewicz identify it with
the relative homotopy group in
degree $r+2$, equivalently $\pi_{r+1}F$.
Deck transformations give the module
action throughout. In particular
\[
 \pi_{r+2}(B\mathcal E)\to\pi_{r+2}Y
 \to\pi_{r+1}F\to\pi_{r+1}(B\mathcal E)
 \to\pi_{r+1}Y\to0
\]
is exact. The module measures more
than an unqualified difference of ranks.

For a counterexample to an always-degree-two
claim take $m\ge3$ and
\[
 Q_m=\{\textstyle\sum_{i=0}^m z_i^2=1\}\subset\C^{m+1},
 \qquad (X,D)=(Q_m\times\C,Q_m\times\{0\}).
\]
Writing $z=x+iy$ gives $x\cdot y=0$ and
$|x|^2=1+|y|^2$. The map
$(x,y)\mapsto(x/\sqrt{1+|y|^2},y)$
identifies $Q_m$ with $TS^m$, so
$X\simeq S^m$. The normal meridian
injects isomorphically into
$\pi_1(Q_m\times\C^*)=\Z$.
The ordinary exit category has a
closed object with trivial group,
an open object with group $\Z$,
and a single crossing arrow.
Its closed object is initial.
Thus its classifying space is
contractible, the comparison fibre
is $S^m$, and its first nonzero
group is $\pi_m=\Z$.
This can occur arbitrarily far
above degree two.
\end{proof}

